\documentclass[10pt,a4paper]{amsart}
\usepackage[margin=19mm]{geometry}
\usepackage{amsmath,amssymb,mathtools}
\usepackage[scr=rsfs,cal=euler]{mathalfa}
\usepackage{xfrac}
\usepackage[unicode,hidelinks,bookmarks=false]{hyperref}
\newcommand{\ie}{i.e.\ }
\newcommand{\cf}{cf.\ }
\newcommand{\resp}{respectively\ }
\newcommand{\Subsection}{\subsection}
\providecommand{\nobreakdash}{\nobreak\hbox{-}}
\setlength{\emergencystretch}{2em}
\multlinegap0pt
\usepackage{xcolor}
\usepackage{amssymb}
\newtheorem{thm}{Theorem}
\newtheorem{prop}{Proposition}
\def\a{\alpha}
\def\bz{{\mathbb Z}}
\def\cb{{\mathcal B}}
\def\cd{{\mathcal D}}
\def\ch{{\mathcal H}}
\def\cs{{\mathcal S}}
\def\cu{{\mathcal U}}
\def\f{\varphi}\def\ff{\phi} \def\F{\Phi}
\def\ga{{\mathfrak A}} \def\gpa{{\mathfrak a}}
\def\om{\omega} \def\Om{\Omega}
\def\r{\rho}\def\R{\mathop{\rm P}}
\def\tr{\mathop{\rm Tr}}
\title{Dominance and equivalence for states on $C^*$-algebras: Quasi-Invariant states}
\author{Ameur Dhahri, Francesco Fidaleo, Chul Ki Ko, Hyun Jae Yoo}
\date{Source: arXiv:2609.00409v1 (CC BY 4.0)}
\begin{document}
\maketitle

\noindent\emph{Source and licence.} Dominance and equivalence for states on $C^*$-algebras: Quasi-Invariant states. Version arXiv:2609.00409v1 (CC BY 4.0), distributed by arXiv under the Creative Commons Attribution 4.0 International licence, http://creativecommons.org/licenses/by/4.0/, as recorded in the arXiv licence metadata for this exact version and shown on the abstract page https://arxiv.org/abs/2609.00409v1. Full licence notice: \texttt{licenses/arxiv-2609.00409-CC-BY-4.0.txt}.\par\smallskip

\noindent\textit{Editorial scope.} Counted unit: the article's prop prop (source0.tex lines 1364--1384). The article's complete original proof of it is delivered with it (source0.tex lines 1385--1428, one \texttt{proof} environment of 44 lines). No mathematics is written, repaired or re-derived: the statement and the proof are the article's own lines, in the article's own notation, and no retained line is substituted or re-flowed.  Retained uncounted as the source-local context the proof needs: lines 501--513; lines 587--591; lines 664--676; lines 708--716; lines 1100--1111; lines 1311--1315; lines 1324--1327; lines 1329--1333; lines 1348--1351; lines 1359--1362.  Nothing was omitted from any retained span, so the package carries no omission record.  No retained line contains a citation, so the wrapper carries no bibliography and no bibliography entry.  No retained line contains a figure environment, an image-inclusion command or drawing code, so no image is delivered and the package declares no asset dependency.  Editorial formatting only, in addition: an AMS \texttt{amsart} preamble that restates the article's own declarations and macros used by the retained lines, namely the environments \texttt{thm}, \texttt{prop} on separate counters, together with the standard packages those lines need.  The retained environments print in this document as Theorem 1, Proposition 1 in order of appearance, whereas the article prints the same items with the numbering of its own full document.  The complete native source of the article as distributed in the arXiv e-print for this exact version is bundled unchanged as \texttt{sources/209023/source0.tex}.
\begin{thm}
\label{main}
Let $\ga$ be a $C^*$-algebra and $\f,\om\in\ga^*_+$. The following assertions are equivalent:
\begin{enumerate} 
\item[(i)] $\om\prec\f$;
\item[(ii)] there exists a positive self-adjoint operator $T_{\om/\f}$, affiliated with $\pi_\f(\ga)'$, such that 
$\pi_\f(\ga)\xi_\f$ is a core for $T_{\om/\f}^{1/2}$ and
$$
\om(x)=\big\langle T_{\om/\f}^{1/2}\pi_\f(x)\xi_\f,T_{\om/\f}^{1/2}\xi_\f\big\rangle\,,\quad x\in\ga\,.
$$
\end{enumerate}
If one of the above equivalent conditions holds true and, in particular, (ii), the operator $T_{\om/\f}$ is uniquely determined.
\end{thm}

\begin{prop}
\label{pi}
Let $\omega,\varphi\in\cs(\ga)$ such that $\omega\prec\varphi$. With $P_{\om/\f}$ denoting the projection onto the closed subspace $[\pi_\f(\ga)T^{1/2}_{\om/\f}\xi_\f]$, $P_{\om/\f}\in\pi_\f(\ga)'$ and 
$(P_{\om/\f}\pi_\f,  P_{\om/\f}\ch_\f, T^{1/2}_{\om/\f}\xi_\f)$ is a GNS representation of $\om$.
\end{prop}

\begin{thm}
\label{invinv}
Let $\ga$ be a $C^*$-algebra and $\f,\om\in\cs(\ga)$ with $\om\prec\f$. 
With the notations of Proposition \ref{pi}, the following are equivalent:
\begin{itemize}
\item[(i)] $T_{\om/\f}^{1/2}\xi_\f$ is cyclic for $\pi_\f(\ga)$;
\item[(ii)] $T_{\om/\f}$ is injective;
\item[(iii)] $P_{\om/\f}=I_{\mathcal H_\varphi}$;
\item[(iv)]  $\f\prec\om$.
\end{itemize}
If one of the above conditions holds true and, in particular (ii),
we have $T_{\f/\om}=T_{\om/\f}^{-1}$.
\end{thm}

\begin{thm}
\label{twogns}
Let $\ga$ be a $C^*$-algebra, $\a$ a $*$-automorphism and $(\f\circ\a)\prec\f\in\cs(\ga)$. Then there exists a (unique) isometry $U_\a$ acting on $\ch_\f$ satisfying $U_\a U_a^*=P_{(\f\circ\a)/\f}$, such that
$$
U_\a\xi_\f=T^{1/2}_{(\f\circ\a)/\f}\xi_\f\,\,\text{and}\,\, U^*_\a\pi_\f U_\a=\pi_\f\circ\a\,.
$$

In addition, $\f\circ\a\sim\f\iff U_\a$ is unitary.
\end{thm}

\begin{thm}[chain rule]
\label{lenire}
Consider an action $G\stackrel{\a_g}{\curvearrowright}\ga$ of a group $G$ on the $C^*$-algebra $\ga$ by $*$-automorphisms $\a_g$. With $D=\pi_\f(\ga)\xi_\f$, 
\begin{equation}
\label{relautz}
A_o:=U_g T_h^{1/2}U_g^*T_g^{1/2}\lceil_D=U_g T_h^{1/2}T_{g^{-1}}^{-1/2}U_g^*\lceil_D
\end{equation}
is closable, and we get
$$
T_{hg}=A_o^*\overline{A_o}\,.
$$
\end{thm}

\begin{equation}
\label{ohrs}
\rho=\sum_{n\in\mathbb Z}\mu_ne_{nn}\,,\quad
\mu_n=\frac{1}{Z}e^{- n^2},\,\, Z=\sum_{n\in\mathbb Z}e^{-n^2}\,,
\end{equation}

\begin{equation}
\label{ohrs1}
g_k(a):={\rm ad}_{V^k}(a)=V^{k}aV^{-k}\,,\quad k\in\mathbb Z,\,\,a\in\ga\,. 
\end{equation}

\begin{equation}
\label{dual}
g_k^{\rm t}(\r)=
V^{-k}\rho V^k=\sum_{n\in\mathbb Z}\mu_{n+k}e_{nn}\,.
\end{equation}

\begin{equation}
\label{norm1}
T^{1/2}_{g_k}\xi_\f=\rho^{1/2}h_k^{1/2}=																\sum_{n\in\mathbb Z}\mu_{n+k}^{1/2}e_{nn}\,.
\end{equation}

\begin{equation}
\label{rhok}
V^{-k}\rho V^k=\rho^{1/2}h_k\rho^{1/2}=(\rho^{1/2}h_k^{1/2})(h_k^{1/2}\rho^{1/2})\,.
\end{equation}

\begin{prop}
For the state $\f$ on $\cb(\ell^2(\bz))$ with density $\r$ given in \eqref{ohrs}, the following properties hold true:
\begin{itemize}
\item[(i)] $\f$ is quasi invariant under the action of $\bz$ given in \eqref{ohrs1};
\item[(ii)] for the unitary implementation $U$ in Theorem \ref{twogns}, we have
\begin{eqnarray*}
U_{g_k}x=V^{-k}xV^k\equiv\pi_{\rm L}(V)^{-k}\pi_{\rm R}(V)^{k}x\,,\quad x\in\mathcal H_\f,\,\,k\in\mathbb Z;
\end{eqnarray*}
\item[(iii)] $U_{g_k}T_{g_l}^{1/2}U_{g_k}^*T_{g_k}^{1/2}$ is essentially selfadjoint on 
$\cd:=\pi_\f(\ga)\xi_\f=\pi_{\rm L}(\ga)\r^{1/2}=\ga\r^{1/2}$, whose closure is $T_{g_{k+l}}^{1/2}$.
\end{itemize}
As a consequence, 
\begin{itemize}
\item[(1)] the chain-rule involving the quasi-invariance of $\f$ (cf. Theorem \ref{lenire}) leads to
$$
T_{g_{k+l}}^{1/2}=U_{g_k}T_{g_l}^{1/2}U_{g_k}^*T_{g_k}^{1/2}
=U_{g_l}T_{g_k}^{1/2}U_{g_l}^*T_{g_l}^{1/2}\,,\quad k,l\in\bz\,;
$$
\item[(2)] $\bz\ni k\mapsto U_{g_k}\in\cu\big(L^2(\cb(\ell^2(\bz)))\big)$ provides a representation of $\bz$.
\end{itemize}
\end{prop}

\begin{proof}
(i) For $a\in\ga$, using repeatedly \eqref{rhok} we get
\begin{align*}
\f\circ g_k(a)=&\langle \pi_\f(g_k(a))\xi_\f,\xi_\f\rangle
              =\langle g_k(a)\rho^{1/2},\rho^{1/2}\rangle\\
							=&\tr(\rho(V^{k}aV^{-k}))
							=\tr((V^{-k}\rho V^{k})a)\\
							=&\tr((a\rho^{1/2}h_k^{1/2})(h_k^{1/2}\rho^{1/2}))
							=\tr((h_k^{1/2}\rho^{1/2})(a\rho^{1/2}h_k^{1/2}))\\
							=&\langle a\rho^{1/2}h_k^{1/2},\rho^{1/2}h_k^{1/2}\rangle
							=\langle \pi_\f(a)T^{1/2}_{g_k}\xi_\f,T^{1/2}_{g_k}\xi_\f\rangle\,.
\end{align*}
This proves that $\f\circ g_k\prec \f$ by Theorem \ref{main}. Since $T_{g_k}$ is injective, by Theorem \ref{invinv}, $\f\circ g_k\sim\f$ and $\f$ is quasi-invariant under the action of $\bz$ given in \eqref{ohrs1}.

(ii) For $x\in\cb(\ell^2(\mathbb Z))\equiv\ga$ and $y\in L^2(\cb(\ell^2(\mathbb Z)))\equiv\ch_\f$, we get
$$
\big(U_{g_k}^*xU_{g_k}\big)y=V^k(x(V^{-k}yV^k))V^{-k}=(V^k xV^{-k})y=g_k(x)y\,.
$$
By taking into account of \eqref{norm1} and \eqref{dual}, we have
$$
T^{1/2}_{g_k}\xi_\f=\rho^{1/2}h_k^{1/2}=V^{-k}\rho^{1/2} V^k=U_{g_k}\xi_\f\,,
$$
and thus the conditions in Theorem \ref{twogns} are satisfied.

(iii) Put $A:=U_{g_k}T_{g_l}^{1/2}U_{g_k}^*T_{g_k}^{1/2}$. Taking into account \eqref{norm1},
we compute
\begin{eqnarray*}
A\pi_\f(a)\xi_\f&=&U_{g_k}T^{1/2}_{g_l}U_{g_k}^*T_{g_k}^{1/2}\pi_\f(a)\xi_\f
                  =U_{g_k}V^ka(V^*)^k\rho^{1/2}h_l^{1/2}\\
									&=&(V^*)^{k}V^ka(V^*)^k\rho^{1/2}h_l^{1/2}V^k
									=a(V^*)^k\rho^{1/2}h_l^{1/2}V^k\\
									&=&a\sum_{n\in\mathbb Z}\mu_{n+l}^{1/2}\;(V^*)^ke_{nn}V^k
									=a\sum_{n\in\mathbb Z}\mu_{n+l}^{1/2}\;e_{(n-k)(n-k)}\\
									&=&a\sum_{n\in\mathbb Z}\mu_{n+l+k}^{1/2}\;e_{nn}
									=T^{1/2}_{g_{k+l}}\pi_\f(a)\xi_\f\,.
\end{eqnarray*}
Since $\pi_\f(a)\xi_\f$ is a core for $T^{1/2}_{g_{k+l}}$, it follows that 
$$
\overline{A\lceil_{\pi_\f(a)\xi_\f}}=\overline{T^{1/2}_{g_{k+l}}\lceil_{\pi_\f(a)\xi_\f}}=T^{1/2}_{g_{k+l}}
$$ 
and (iii) holds true.

To end the proof, notice that (1) and (2) are direct consequences of (i)-(iii).
\end{proof}

\end{document}
